\section{The question}
\label{sec:question}
An exhaustive search at order 32 spends almost all its effort in prefixes that pass every local test yet have
no completion. At mid-depth this is nearly universal: of the prefixes that survive the native checks and
forward checking, 90--100\% are dead (Section~\ref{sec:local}). The question of this paper is whether those
deaths have \emph{reasons shorter than the search that discovers them}: rules, certificates or relaxations
that refute a dead prefix with far less work than exhausting its subtree, and that recur across positions
and orders. Each candidate was tested against a gate fixed in advance, against negative controls, and,
for the strongest results, by an independent implementation.

\section{Method}
\label{sec:method}
\paragraph{Ground truth.} At orders 12 to 15 every Costas array is listed by the project's exact search, so a
prefix is completable exactly when some array contains it. At larger orders completability is decided by
a SAT solver and, for subtree sizes, by the native search; no relaxation is ever used as ground truth.

\paragraph{Prefixes.} A prefix fills columns from the left. It is \emph{locally admissible} if its rows are
distinct, no displacement repeats, and every unplaced column still has a legal row (forward checking).

\paragraph{Sampling.} Prefixes are drawn by random walks through the locally admissible tree. A walk reaches
a prefix with probability $\prod_i 1/b_i$, where $b_i$ is the number of legal choices at step $i$, so weighting
by $\prod_i b_i$ gives estimates for the uniform depth layer; weighting dead prefixes by their native subtree
size gives estimates of the share of \emph{work} affected. Weighted estimates are reported with their
effective sample sizes $N_{\mathrm{eff}}=(\sum w)^2/\sum w^2$.

\section{Local obstructions}
\label{sec:local}
\paragraph{Minimal obstructions.} For dead order-12 prefixes at depth 5, deleting dots one at a time while the
set stays infeasible gives minimal infeasible dot sets: sets of dots with no repeated displacement that lie
in no Costas array. They are tiny, almost always four dots, but scattered: 60 obstructions came in 56
shapes, and none of twelve shapes tested is infeasible at every position at order 12 or 13.

\paragraph{Negative control.} The same pipeline was run against 7{,}852 random permutations, the number of
order-12 Costas arrays, on the same prefixes:

\begin{table}[H]
\centering
\small
\begin{tabular}{@{}lrrrr@{}}
\toprule
depth 5 & dead & 3-dot obstructions & 4-dot & recurring shapes \\
\midrule
Costas arrays & $86\%$ & 19 & 319 & $29\%$ \\
random permutations & $91\%$ & 6 & 352 & $22\%$ \\
\bottomrule
\end{tabular}
\caption{Minimal obstructions against the Costas arrays and against an equally sparse random set, on the same
400 sampled prefixes.}
\label{tab:control}
\end{table}

\noindent The bulk of the obstruction statistics is reproduced by sparsity alone. The Costas-specific residue
is small: more three-dot obstructions, and somewhat more recurrence.

\paragraph{The residue is positional.} Exhaustively at order 12, 3{,}508 of the 288{,}600 three-dot sets with no
repeated displacement lie in no Costas array: $1.2\%$, against about $0.27\%$ expected for a random set of the
same size. The excess concentrates at corners: $0.84\%$ of triples with no corner dot are infeasible, $4.9\%$
with one, and $62\%$ with two. No three-dot shape is infeasible at every position at both orders 12 and 13.

\paragraph{A positive control: corner reduction.} Deleting a corner dot translates the other dots, so order-$n$
arrays with a dot at $(0,0)$ correspond exactly to order-$(n-1)$ arrays that admit a corner addition. This
predicts the feasibility of every corner triple at order $n$ from the catalogue at order $n-1$:

\begin{table}[H]
\centering
\small
\begin{tabular}{@{}rrrr@{}}
\toprule
order & arrays with a dot at $(0,0)$ & admissible arrays, order $n-1$ & triples predicted correctly \\
\midrule
13 & 509 & 509 & 8{,}676 of 8{,}676 \\
14 & 734 & 734 & 12{,}132 of 12{,}132 \\
15 & 616 & 616 & 16{,}513 of 16{,}513 \\
\bottomrule
\end{tabular}
\caption{Corner reduction as a lifted obstruction: exact at three orders, with no refitting.}
\label{tab:corner}
\end{table}

\noindent Plain infeasibility one order lower explains none of the 1{,}599 infeasible corner triples; the whole
effect is the small size of the corner-admissible subclass. The apparatus recovers real cross-order
structure, and the structure it recovers is the corner surgery of Paper~2.

\section{Clause learning, matching and cutting planes}
\label{sec:proofs}
\paragraph{Incremental clause learning.} Solving 36 order-16 sibling cubes with one persistent CDCL solver
per family, instead of a fresh solver per cube, saved $8$--$18\%$ of conflicts; related siblings benefited
slightly more than unrelated cubes.

\paragraph{The native/SAT frontier.} For dead prefixes that survive all local checks, the native search's
subtree size was compared with the conflicts a fresh CDCL solver needs to refute the prefix. At order 14 the
median ratio was $1.4$--$1.6$ native nodes per conflict, at order 16 $1.16$, and no prefix showed even a
five-fold advantage: CDCL finds no shorter reasons than the search.

\paragraph{Hall's condition.} A matching test of unplaced columns against unused rows catches $0$--$2\%$ of dead
prefixes at mid-depth, rising to $62\%$ only at depth 9 of 12, where little work remains.

\paragraph{Cutting planes.} Dead prefixes below the LP transition of Section~\ref{sec:lp} were refuted by Exact,
a cutting-planes pseudo-Boolean solver descended from RoundingSat~\cite{elffers2018}, on the native
integer model of that section. Every prefix was refuted and every live control satisfied, but the effort
tracked the native subtree: medians of $1.3$, $1.2$ and $0.9$ native nodes per conflict at orders 12, 13 and
14, falling as subtrees grow. The build used had no LP integration, and its proofs were not logged.

\section{The pair-marginal relaxation}
\label{sec:lp}
\paragraph{Formulation.} Placement variables $x_{i,r}\in[0,1]$ satisfy the row and column equalities. For each
pair of columns $i<j$, variables $y_{ijrs}\in[0,1]$ over distinct rows carry one unit of mass, with marginals
$\sum_s y_{ijrs}=x_{i,r}$ and $\sum_r y_{ijrs}=x_{j,s}$: the first level of the Sherali--Adams
hierarchy~\cite{sherali1990}. Each displacement class $(d,e)$ has capacity $\sum y\le1$ over the pairs at
lag $d$ and rows differing by $e$. A prefix fixes its placements to 1. Every constraint holds for every
Costas array, so the relaxation is sound at every order. A weaker form replaces the marginals by the
McCormick inequalities $y\le x_{i,r}$, $y\le x_{j,s}$, $y\ge x_{i,r}+x_{j,s}-1$.

\paragraph{Soundness, checked.} No prefix of a real Costas array was cut (60 of 60, at depths 3 to 10 of order
12), nor any of 483 completable prefixes in the confirmation below.

\paragraph{A sharp transition.} Across orders 12 to 16, with 60 to 160 dead prefixes per cell:

\begin{table}[H]
\centering
\small
\begin{tabular}{@{}rrrrrr@{}}
\toprule
order & depth & $k/n$ & weak form & strong form & strong, uniform layer \\
\midrule
12 & 6 & $0.500$ & $5\%$ & $28\%$ & $25\%$ \\
12 & 7 & $0.583$ & $53\%$ & $100\%$ & $100\%$ \\
13 & 7 & $0.538$ & $14\%$ & $74\%$ & $72\%$ \\
13 & 8 & $0.615$ & $75\%$ & $100\%$ & $100\%$ \\
14 & 8 & $0.571$ & $41\%$ & $98\%$ & $99\%$ \\
15 & 8 & $0.533$ & $8\%$ & $63\%$ & $68\%$ \\
15 & 9 & $0.600$ & $57\%$ & $100\%$ & $100\%$ \\
16 & 9 & $0.562$ & $30\%$ & $93\%$ & $94\%$ \\
16 & 10 & $0.625$ & $83\%$ & $100\%$ & $100\%$ \\
\bottomrule
\end{tabular}
\caption{Share of dead prefixes the relaxation proves infeasible, near the transition. Work-weighted rates
agree beyond the transition and are slightly lower inside it.}
\label{tab:transition}
\end{table}

\noindent Logistic fits put the 90\% point $\alpha_{90}$ between $0.512$ and $0.557$ of the order at all five
orders, and a fixed-fraction model fits far better than a fixed number of remaining columns (negative
log-likelihood 314 against 606). Beyond the transition no miss occurred among 1{,}140 sampled dead
prefixes, an approximate 95\% upper bound of $0.26\%$ on the miss rate, with the weighted estimates agreeing
at effective sample sizes of 32 to 146 per cell. An independent implementation, with its own enumerator,
generator and sampler, reproduced the transition and the weak--strong contrast at orders 12 to 14.

\paragraph{The certificates are global.} Shrinking the displacement capacities to a minimal set of lags that
still refutes each prefix leaves 3 to 5 lags, but the sets vary: 17 distinct among 20 prefixes at order 12.
A per-lag matching test, Hall's condition on column pairs and differences, explains only 5 of 25 and 1 of 16
of the certified prefixes. The refutations couple several lags through shared columns.

\section{Order-32 geometry}
\label{sec:geometry}
The validated estimator of Paper~9, applied per depth to the order-32 tree with all eight symmetries,
locates the transition at depths 17--18 and the tree's widest level at depth 20. Calibrated against exact
counts at orders 12 to 14, the per-depth estimates were within $0.3\%$.

\begin{table}[H]
\centering
\small
\begin{tabular}{@{}rrrr@{}}
\toprule
cut depth & locally admissible prefixes & share of the tree at or above it & best possible saving \\
\midrule
16 & $1.91\times10^{18}$ & $3.6\%$ & $28\times$ \\
17 & $5.00\times10^{18}$ & $11.4\%$ & $8.7\times$ \\
18 & $8.20\times10^{18}$ & $27.9\%$ & $3.6\times$ \\
19 & $7.41\times10^{18}$ & $52.0\%$ & $1.9\times$ \\
\bottomrule
\end{tabular}
\caption{The order-32 frontier, and the most any rule complete at a given depth could save if it cost
nothing.}
\label{tab:frontier}
\end{table}

\noindent Solving the relaxation once per frontier prefix is out of the question. More fundamentally, even a
free, perfect replica of its verdict at depth 17 saves at most a factor of about nine. The depth of the transition was measured at orders 12 to 16, near the
peak of the Costas counts; that it stays near the same fraction of the order as the counts fall, which
they do exponentially fast in proportion to all permutations~\cite{warnke2023density}, was checked at orders
18 and 20, beyond the peak, for binary displacement consistency (Section~\ref{sec:late}): there too it
refutes almost nothing at half depth and about three quarters of prefixes at $0.55n$. That the fraction
stays put through the far sparser orders 24 to 32 remains an assumption.

\section{Triple consistency}
\label{sec:triple}
The next level of the hierarchy adds, for every triple of columns, a distribution over distinct row triples
whose three projections equal the three pair distributions. An independent order-12 implementation, on
exactly equivalent prefix-reduced models, found it strictly stronger: it certified all 12 dead prefixes
sampled at depth 6, where the pair relaxation certifies about half. At depth 5 ($0.42n$) it certified 2 of 22,
far below the predeclared gate of 90\%. At depth 4 no instance was certified within the time limit. No live
control was cut. The transition moves by about one column, from near $0.58n$ to near $0.50n$ at order 12; at
order 32 that corresponds to depth 16, where the ceiling of Table~\ref{tab:frontier} is $28\times$, for a
relaxation that could not be run at that order.

\section{Late mechanisms}
\label{sec:late}
Three further mechanisms that work on exact displacement occupancy were measured, one by the project's
search and two by an independent implementation within the project, whose code and data are kept in
\texttt{mining/independent/}.

\paragraph{Binary displacement consistency.} At every node, each pair of unplaced columns must admit a pair
of rows that repeats no used displacement, either between the two new dots or against the placed dots,
and rows without such a partner are removed until nothing changes. On complete searches it is sound (every
order returns exactly the known arrays) and strong:

\begin{table}[H]
\centering
\small
\begin{tabular}{@{}rrrrr@{}}
\toprule
order & no lookahead & forward checking & binary consistency & reduction \\
\midrule
10 & 152{,}100 & 73{,}778 & 19{,}715 & $7.7\times$ \\
12 & 3{,}488{,}680 & 1{,}583{,}032 & 284{,}922 & $12.2\times$ \\
13 & 18{,}290{,}043 & 8{,}116{,}650 & 1{,}362{,}891 & $13.4\times$ \\
14 & 83{,}176{,}084 & 36{,}110{,}650 & 5{,}314{,}892 & $15.6\times$ \\
15 & 435{,}090{,}280 & 185{,}553{,}661 & 26{,}624{,}065 & $16.3\times$ \\
\bottomrule
\end{tabular}
\caption{Nodes of the left-to-right search with three levels of lookahead.}
\label{tab:binac}
\end{table}

\noindent The reduction grows with the order, but more slowly at each step, which is consistent with a large
constant factor rather than a lower growth rate. Where it acts settles the question: at orders 13 and 14,
93\% and 97\% of its cuts fall between $0.54n$ and $0.64n$, almost none before $0.5n$, and the smallest
remaining domain at a cut is typically one or two rows. Binary consistency is therefore strong in the
same band as the pair-marginal relaxation, where the tree is widest, and not earlier. Recomputed from
scratch at every node, as here, it is about eight times slower than the plain search at order 15.

Whether it can pay for itself was then tested against a stronger baseline, the search with forward
checking, using one engine and three policies whose solution sets were checked to be identical. Against
forward checking, binary consistency saves only about four times the nodes at order 14, and each call
costs about forty forward-checking nodes. Triggering it only where it acts, from about $0.6n$ or when a
domain is down to one row, still left the search 1.3 to 1.5 times slower at orders 13 and 14. More
decisively, the time spent outside the consistency checks bounds what any implementation could achieve:
even checks that cost nothing would make the search at most 2.0 times faster at order 13 and 1.6 times at
order 14. That ceiling falls with the order. The tested implementation therefore fails a predeclared
production gate of a 1.25-fold speed-up, and an incremental one would have to be about four times cheaper
merely to reach that gate at order 14, against a shrinking margin; binary consistency was not retained in
the production solver. The comparison used only the vertical-flip symmetry, identically for all three
policies, and was run at orders 12 to 14. The two left-to-right policies did not commit forced columns out
of order, whereas the policy with dynamic column choice does so implicitly by always choosing the column
with fewest rows; that asymmetry favours dynamic choice, which nevertheless did not improve elapsed time.
Binary consistency thus remains informative about where pruning power lies, but not, on this evidence, a
route to a faster search.

The same transition holds beyond the peak of the Costas counts. The independent implementation applied
binary consistency to 24 random prefixes per depth at orders 18 and 20, which have 15{,}096 and 6{,}464
arrays. At half depth it refuted 1 prefix of 24 at each order; at depth 10 of 18 it refuted 18 of 24, and at
depth 11 of 20 it refuted 16 of the 21 that forward checking had not already closed. In both cases the
smallest remaining domain at that depth was typically two rows. Among prefixes that survive binary
consistency near half depth, the pair-marginal relaxation still proved some infeasible (2 of 7 at order 18
depth 9, 1 of 8 at order 20 depth 10), so on some prefixes it is strictly the stronger of the two.

\paragraph{Chain-filtered Hall analysis.} The independent implementation kept, for every pair of columns,
only the row pairs supported through every intermediate column, and then tested Hall's condition on each
lag's positions and available differences. Hall deficiencies appear in none of 20 dead prefixes at order 14
depth 6, in 4 at depth 7 and in 19 at depth 8 ($0.57n$); at order 13 depth 7 they appear in 10 of 20,
where the pair-marginal relaxation certifies about three quarters. Splitting on sets with one unit of slack
increased total work for every one of the 49 prefixes tested.

\paragraph{Where completion is decided.} The same implementation counted, at orders 12 to 14, which locally
valid prefixes extend to a Costas array. At depth 3, 96--98\% do; at the midpoint depth $\lceil0.53n\rceil$,
0.8\%, 0.5\% and 0.09\% do. No sound rule can prune much at depth 3 of these orders. The work between the two
depths is nearly uniform across depth-3 prefixes (Gini coefficient about 0.09), and only 30--38\% of it lies
above the midpoint level, as a branching factor near four implies, so compiling that stretch of the search
could save at most a factor of about 1.3. These orders sit near the peak of the Costas counts, which caps
what can be learned from them about early pruning at order 32.

\section{Aggregate relaxations}
\label{sec:aggregate}
At lag $d$ the Costas condition asks for $n-d$ distinct values in a range of $2n-2$: injectivity at density
at most one half. Relaxations that keep only aggregates of the differences, such as sums, sums of squares
or residues, lose which values are occupied, which is where the constraint bites. A direct test: for each
lag, the sum of squares of the differences must lie between the least and greatest values achievable by
$n-d$ distinct allowed differences with the same sum. At order 32, $95.9\%$ of 20{,}000 random permutations,
none of them Costas, satisfy all 31 of these conditions, as do $N_{32}$ and the witness of Paper~7; at
order 12 the share is $62.3\%$, so the conditions weaken as the order grows. Correll and Swanson make the
corresponding observation for a single statistic: the number of positive first differences, which they
show lies within 2.5 of $(n-1)/2$ for every known Costas array, by itself yields only an $O(n^{-1/2})$
density~\cite{correll2023difference}, against the exponential decay of the true
density~\cite{warnke2023density}.

Classifying midpoint states by short-lag signatures falls to the same effect. With a signature recording,
for lags 1 to 4, how many differences are positive and how many are small, signature classes with at least
100 members and no completable state cover $9.35\%$ of the 987{,}664 locally valid order-12 midpoint states,
as the independent implementation found and a recomputation here reproduces exactly. With the completable
labels shuffled at random, such classes cover $8.3$--$9.0\%$: almost all of the apparent exclusion is what
sparsity alone produces.

\section{Implications}
\label{sec:implications}
Three things should be distinguished.
\begin{itemize}
\item \emph{An obstruction language was found.} The pair-marginal relaxation proves dead prefixes infeasible
with linear certificates that no local rule and no resolution-based solver produces, and it does so
completely, in sample, beyond about half depth. That is new, and it is the first strong global obstruction
system for partial Costas permutations that we know of.
\item \emph{A cheap propagator was not.} Every reasoning tested below the transition, whether local
rules, clause learning, cutting planes or triple consistency, either fails to refute early or needs effort
comparable to the native search. The mechanisms that work on exact occupancy, namely the pair-marginal
relaxation, binary displacement consistency and chain-filtered Hall analysis, all become strong in the
same band near half depth, and aggregate statistics carry little beyond sparsity.
\item \emph{Order 32 is not made feasible.} The transition falls where the order-32 tree is widest, which caps
any rule acting there at about a nine-fold saving, or 28-fold for the triple level.
\end{itemize}
Two questions remain open. Why the transition sits near half depth is unexplained; the number of pairs
between placed and unplaced columns, $k(n-k)$, peaks at $k=n/2$, which suggests a mechanism but proves
nothing. And the certificates themselves, reconstructed exactly from dual multipliers, have not yet been
read for a common form.
